\documentclass[10pt]{article}
\usepackage[T1]{fontenc}
\usepackage{lmodern,amsmath,amssymb,amsthm,mathtools,booktabs,microtype}
\usepackage[margin=0.86in]{geometry}
\usepackage[hidelinks]{hyperref}
\usepackage{enumitem}
\setlist{itemsep=2pt,topsep=4pt,parsep=0pt}
\setlength{\parindent}{1.2em}
\setlength{\parskip}{1pt}
\newtheorem{theorem}{Theorem}
\newtheorem{lemma}[theorem]{Lemma}
\newtheorem{proposition}[theorem]{Proposition}
\DeclareMathOperator{\Tr}{Tr}
\DeclareMathOperator{\rank}{rank}
\DeclareMathOperator{\vecop}{vec}
\newcommand{\HH}{\mathcal H}
\newcommand{\Aux}{\mathcal A}
\newcommand{\QQ}{\mathbb Q}
\newcommand{\CC}{\mathbb C}
\newcommand{\al}{\alpha_0}
\newcommand{\eps}{\varepsilon}
\begin{document}
\begin{center}
{\Large\bfseries An exact three-copy undistillability bound}\\[3pt]
{\large Uniform in local dimension, at $\alpha=462713/1000000$}
\end{center}
\begin{abstract}
We give an exact-arithmetic, computer-assisted proof that Werner states with
parameter $0\le\alpha\le0.462713$ are three-copy undistillable in every
finite local dimension. The proof combines a dimension-independent quadratic
inequality for arbitrary complex rank-two matrices with a rational
four-replica certificate. Positivity of 238 strictly positive leading blocks
implies positivity in all sufficiently large dimensions; isometric embedding
then covers every finite dimension. The analytic inequality is proved here.
The accompanying rational data and independent checker comprise the finite
computational part of the proof.
\end{abstract}

\section{Statement and conventions}
Let $F(x\otimes y)=y\otimes x$ on $\CC^d\otimes\CC^d$, and define
\begin{equation}\label{eq:werner}
 \rho_{d,\alpha}=\frac{I-\alpha F}{d(d-\alpha)},\qquad
 \al=\frac{462713}{1000000}=0.462713.
\end{equation}
All partial traces below are unnormalized. For a matrix $X$ on
$\HH=\HH_1\otimes\HH_2\otimes\HH_3$, write
\begin{equation}\label{eq:Q}
 Q_\alpha(X)=\sum_{S\subseteq\{1,2,3\}}(-\alpha)^{|S|}
                      \|\Tr_S X\|_F^2
 =N-\alpha B_1+\alpha^2B_2-\alpha^3T,
\end{equation}
where $N=\|X\|_F^2$, $B_j=\sum_{|S|=j}\|\Tr_S X\|_F^2$,
and $T=|\Tr X|^2$. A missing second singular value is taken to be zero.

\begin{theorem}\label{thm:main}
For arbitrary finite-dimensional complex spaces $\HH_1,\HH_2,\HH_3$ and
every matrix $X$ on their tensor product with $\rank X\le2$,
\[
 Q_{\al}(X)\ge0.
\]
Consequently, for every integer $d\ge2$ and $0\le\alpha\le\al$,
$\rho_{d,\alpha}$ is three-copy undistillable. It has negative partial
transpose when $\alpha>1/d$; in particular the theorem gives a nonempty NPT
interval $1/d<\alpha\le\al$ for every $d\ge3$.
\end{theorem}

Here three-copy undistillability is the Schmidt-rank-two condition
\[
 \langle\psi|(\rho_{d,\alpha}^{\Gamma})^{\otimes3}|\psi\rangle\ge0
 \quad\text{for every $\psi$ of Schmidt rank at most two,}
\]
with the three copies grouped by party. Indeed
$F^\Gamma=|\Omega_d\rangle\langle\Omega_d|$ for
$\Omega_d=\sum_{a=1}^d e_a\otimes e_a$, and vectorization identifies
Schmidt rank with matrix rank and the expectation with
$Q_\alpha(X)/[d(d-\alpha)]^3$. A Schmidt-rank-two vector is the pullback
of a two-qubit vector under local filters, giving the usual two-qubit-output
criterion for distillation. The eigenvalue $1-d\alpha$ of
$I-\alpha|\Omega_d\rangle\langle\Omega_d|$ gives the asserted NPT condition.
The theorem concerns three physical copies, not arbitrary copy number.

\section{A quadratic inequality and its determinant}
The analytic ingredient is the following inequality, proved in full in
Appendix~\ref{app:analytic}:
\begin{equation}\label{eq:A}
 \Aux(X):=\frac{17}{8}N-B_1+\frac12B_2-\frac5{16}T
 \ \ge\ \frac{(\sigma_1-\sigma_2)^2}{4}\ge0
 \qquad(\rank X\le2).
\end{equation}
It applies to arbitrary complex $X$ and unequal factor dimensions.

For four vectors $r=(u_0,v_0,u_1,v_1)$ in $\HH$, put
$D_j=u_jv_j^\dagger$. For real subset weights $w=(w_S)$ define
\[
 (H_w)_{ij}=\sum_S w_S\langle\Tr_S D_i,\Tr_S D_j\rangle_F.
\]
Then $c^\dagger H_wc=\sum_Sw_S\|\Tr_S(c_0D_0+c_1D_1)\|_F^2$.
Every combination $c_0D_0+c_1D_1$ has rank at most two. Thus \eqref{eq:A} implies
$H_a\succeq0$ and $\det H_a\ge0$ for the weights
\begin{equation}\label{eq:aweights}
 a_\varnothing=17/8,\quad a_S=-1\ (|S|=1),\quad
 a_S=1/2\ (|S|=2),\quad a_{123}=-5/16.
\end{equation}

We make the polynomial determinant explicit. Replicas are ordered
$(u_0,v_0,u_1,v_1)$ and numbered $0,1,2,3$. A permutation $p\in S_4$
maps input position $j$ to output position $p(j)$; products act right to left.
For $e,f\in\{0,1\}$ set
\[
 d_{ef}=(01)^e(23)^f,\qquad
 c_{ef}=(03)^e(12)^f(02)(13).
\]
Writing $e_i=1_{i\in S}$ and $f_i=1_{i\in R}$, define the formal tensor
\begin{equation}\label{eq:Kw}
 K(w)=\sum_{S,R\subseteq\{1,2,3\}}w_Sw_R
 \left(\bigotimes_{i=1}^3 P_{d_{e_if_i}}
       -\bigotimes_{i=1}^3 P_{c_{e_if_i}}\right).
\end{equation}
Regrouping the four replicas by physical site, contraction of indices gives
\begin{equation}\label{eq:det}
 \langle u_0\otimes v_0\otimes u_1\otimes v_1,
 K(w)(u_0\otimes v_0\otimes u_1\otimes v_1)\rangle=\det H_w.
\end{equation}
The first permutation product gives $(H_w)_{00}(H_w)_{11}$ and the
second gives $(H_w)_{01}(H_w)_{10}$. Let $K_\alpha=K(w)$ with
$w_S=(-\alpha)^{|S|}$.

Let $\mathcal S$ denote averaging over physical-site permutations,
simultaneous conjugation of replica labels by
\[
 V_4=\{1,(01)(23),(02)(13),(03)(12)\},
\]
and adjoint. Put $C=\mathcal S K(a)$. These operations preserve
four-vector products and nonnegativity of their real expectations, so
\begin{equation}\label{eq:Cpositive}
 \langle r,Cr\rangle\ge0
 \quad\text{for every four-vector product $r$.}
\end{equation}
This does not assert positivity of $C$ on arbitrary entangled vectors of the
four-replica space.

\section{The rational coefficient certificate}
Use the eight transpose sets
\[
 \mathcal T=\{\varnothing,\{0\},\{1\},\{2\},\{3\},
                         \{0,1\},\{0,2\},\{0,3\}\}.
\]
The supplied rational arrays define
$Y_\tau=\sum_{\mathbf p\in S_4^3}y_{\tau,\mathbf p}
P_{p_1}\otimes P_{p_2}\otimes P_{p_3}$ and a tensor $E$ satisfying
the exact coefficient identity
\begin{equation}\label{eq:cert}
 K_{\al}=\sum_{\tau\in\mathcal T}Y_\tau+E+\mu C,
 \qquad
 \mu=\frac{1581841004766715}{18014398509481984}>0.
\end{equation}
All coefficients are rational constants independent of dimension.
The tensor $E$ is supported only on monomials for which one of
\[
 p_i(0)=2,\quad p_i(2)=0,\quad p_i(1)=3,\quad p_i(3)=1
\]
holds for all three sites $i$, with the same choice throughout. Each such
monomial contains one of the full-vector inner products
$\langle u_1,u_0\rangle$, $\langle u_0,u_1\rangle$,
$\langle v_1,v_0\rangle$, $\langle v_0,v_1\rangle$ as a factor.
It therefore vanishes on orthogonal left and right singular frames.
There are 819 such monomials, forming 40 symmetry orbits.

Write $\Gamma_\tau$ for transpose on the full replica factors in $\tau$,
at all three physical sites. If
\begin{equation}\label{eq:familypositive}
 \Gamma_\tau(Y_\tau)\succeq0 \quad(\tau\in\mathcal T),
\end{equation}
then $\langle r,Y_\tau r\rangle
=\langle r^{*\tau},\Gamma_\tau(Y_\tau)r^{*\tau}\rangle\ge0$;
the selected full vectors are conjugated in $r^{*\tau}$.
Equations~\eqref{eq:Cpositive}--\eqref{eq:familypositive} imply
$\det H_{w(\al)}\ge0$ on every orthogonal singular frame.

The diagonal entries are positive for nonzero dyads: with $F_i$ denoting
the local swaps between two replicas,
\begin{equation}\label{eq:rankone}
 Q_\alpha(uv^\dagger)
 =\langle u\otimes v,\prod_{i=1}^3(I-\alpha F_i)(u\otimes v)\rangle
 \ge(1-\alpha)^3\|u\|^2\|v\|^2\qquad(0\le\alpha<1).
\end{equation}
The swaps commute and each factor is bounded below by $(1-\alpha)I$.
Thus the determinant inequality makes this two-by-two Hermitian restriction
positive semidefinite. Applying it to the SVD of an arbitrary rank-two $X$
proves $Q_{\al}(X)\ge0$ whenever \eqref{eq:familypositive} holds.

\section{Certification in every dimension}
\subsection*{Complete polynomial representations}
At one physical site, the 24 operators
$D_{\tau,p}=\Gamma_\tau(P_p)$ have pairing diagrams with four input and
four output vertices. For $P_p$, input $j$ is paired with output $p(j)$;
transpose exchanges the input and output vertices in $\tau$. A matrix entry
is the product of the Kronecker deltas along the four edges. Composition is
concatenation of diagrams, with one factor $d$ per closed loop. Adjoint
reflects the diagram; trace closes matching input/output vertices.

For $d\ge4$, these 24 operators are linearly independent: assign four
distinct labels to the edges of any selected diagram. This matrix entry
is one for that diagram and zero for every other pairing. Their physical
algebra therefore has dimension 24.

The second data file supplies matrices $\rho_{\tau,r}(p;d)$ and symmetric
Grams $G_{\tau,r}(d)$ over $\QQ[d]$, and polynomial multiplicities
$m_{\tau,r}(d)$. The checker verifies the following as polynomial identities:
\begin{align}
 \rho_r(p)\rho_r(q)&=d^{\ell(p,q)}\rho_r(p\cdot q),
 &\rho_r(p)^T G_r&=G_r\rho_r(p^*),\label{eq:algebra}\\
 \Tr_{\rm phys}D_{\tau,p}
 &=\sum_r m_{\tau,r}(d)\operatorname{tr}\rho_{\tau,r}(p;d),
 &\sum_r(\dim\rho_{\tau,r})^2&=24.\label{eq:traces}
\end{align}
Here $p\cdot q$ labels the concatenated diagram and $\ell(p,q)$ counts
its closed loops; for a nonempty transpose set this is not generally
ordinary permutation multiplication.
It also checks $\rho_r(1)=I$, $G_r(d)\succ0$ and $m_r(d)>0$ for every
real $d\ge4$. Positivity follows by substituting $d=4+u$: every Gram
leading principal minor and multiplicity has nonnegative rational
coefficients and a positive constant term.

These checks prove completeness. The map
$\Phi(a)=\bigoplus_r\rho_r(a)$ is a homomorphism. If $\Phi(a)=0$,
\eqref{eq:traces} applied to $a^\dagger a$ gives
$\Tr_{\rm phys}(a^\dagger a)=0$, hence $a=0$. Domain and codomain both
have dimension 24, so $\Phi$ is an isomorphism. Conjugating each block by
$G_r^{1/2}$ makes it a $*$-isomorphism. Physical positivity is therefore
equivalent to positivity in every supplied block. No irreducible component
is left unchecked.

\subsection*{Strict leading blocks}
For each block, let $k_r$ be the largest degree in $d$ of any entry of
any $\rho_r(p;d)$, and let $g_r$ be the largest degree of a Gram entry.
With $t=1/d$, set
\[
 R_{r,p}(t)=t^{k_r}\rho_r(p;1/t),\qquad
 \overline G_r(t)=t^{g_r}G_r(1/t).
\]
Both are polynomial; the checker verifies
$\overline G_r(0)\succ0$. The complete degree data are:
\begin{center}\small
\begin{tabular}{@{}llllr@{}}\toprule
Transpose representative & Block orders & $k_r$ & $g_r$ & Triple blocks\\\midrule
$\varnothing$ & $1,1,2,3,3$ & $0,0,0,0,0$ & $0,0,0,0,0$ & 35\\
$\{0\}$ & $1,1,2,3,3$ & $0,0,0,1,1$ & $0,0,0,1,1$ & 35\\
Each of $\{0,1\},\{0,2\},\{0,3\}$ & $1,1,1,1,2,4$ & $0,0,0,0,2,1$ & $0,0,0,0,1,1$ & 56\\\bottomrule
\end{tabular}
\end{center}
For a triple $\mathbf r=(r_1,r_2,r_3)$ define, using the same raw coefficients
as in \eqref{eq:cert},
\begin{align}
 M_{\tau,\mathbf r}(t)
 &=\sum_{p_1,p_2,p_3}y_{\tau,\mathbf p}
                  \bigotimes_{j=1}^3R_{\tau,r_j,p_j}(t),\label{eq:normalized}\\
 B_{\tau,\mathbf r}(t)
 &=\left(\bigotimes_j\overline G_{\tau,r_j}(t)\right)
                         M_{\tau,\mathbf r}(t).
\end{align}
These are polynomial, and $B$ is symmetric by \eqref{eq:algebra} and
the verified Hermitian coefficient symmetries. At positive $t=1/d$, $B$
is the actual Gram-weighted physical block of $\Gamma_\tau(Y_\tau)$
multiplied by the positive scalar $t^{\sum_j(k_{r_j}+g_{r_j})}$.
This changes only how a block is tested, not its coefficients or positivity.

For every representative triple, the exact certificate satisfies
\begin{equation}\label{eq:margin}
 G_0(M_0-\eps I)\succeq0,
 \qquad \eps=\frac1{2000000000},\qquad
 G_0=\bigotimes_j\overline G_{r_j}(0),\quad M_0=M_{\tau,\mathbf r}(0).
\end{equation}
Here $I$ is the matrix identity, not the leading image of the identity
diagram, which may vanish. Since $G_0\succ0$, \eqref{eq:margin} implies
$B(0)\succeq\eps G_0\succ0$.

The coefficient symmetries identify reordered triples of block labels;
$V_4$ identifies the four singleton transpose sets. Complementary transpose
sets differ by full transpose, preserving positivity. Consequently the
complete list comprises
$\binom73+\binom73+3\binom83=238$ representative blocks, of order at most 64.
There are finitely many continuous matrices $B(t)$, all strictly positive
at zero. They are therefore simultaneously positive for all sufficiently
small positive $t$. Thus \eqref{eq:familypositive} holds for every
sufficiently large integer $d$.

\subsection*{Descent and the parameter interval}
Given arbitrary factor dimensions $d_i$, choose a sufficiently large common
dimension $D\ge\max_i d_i$ and isometries $V_i:\HH_i\to\CC^D$.
For $\widetilde X=(\bigotimes_iV_i)X(\bigotimes_iV_i)^\dagger$,
\[
 \Tr_S\widetilde X=(\bigotimes_{i\notin S}V_i)
                  (\Tr_S X)(\bigotimes_{i\notin S}V_i)^\dagger.
\]
Rank and all partial-trace Frobenius norms are preserved. Hence
$Q_{\al}(X)=Q_{\al}(\widetilde X)\ge0$. An explicit dimension cutoff is
unnecessary: an arbitrarily large certified dimension can be chosen for
each proposed finite-dimensional matrix.

For completeness the lower-parameter interval follows without any
dimension-dependent estimate. Put
$A_\alpha=I-\alpha|\Omega_d\rangle\langle\Omega_d|$.
Schmidt-rank-two positivity of $A_{\al}^{\otimes3}$ implies that of
$A_{\al}^{\otimes k}$ for $k\le3$, by adjoining product vectors with
positive spectator expectations. For $0\le\alpha\le\al$,
\[
 A_\alpha=(\alpha/\al)A_{\al}+(1-\alpha/\al)I.
\]
Expanding its third tensor power gives a nonnegative combination of
$A_{\al}^{\otimes k}\otimes I$ on subsets of the copies. The expectation
of each term on a Schmidt-rank-two vector is a sum of nonnegative
expectations, obtained by resolving the spectator systems in local product
bases; each resulting vector still has Schmidt rank at most two.
This proves the entire interval in Theorem~\ref{thm:main}.

\section{Finite exact checks and the accompanying data}
The computational proof uses two data files, treated as untrusted input:
\begin{itemize}
\item \texttt{certificate/certificate.json.gz}: the eight rational
coefficient tensors, ideal coefficients, $\mu$, $\al$, and $\eps$;
\item \texttt{certificate/representations.json.gz}: the polynomial matrices,
Grams and multiplicities defining \eqref{eq:algebra}--\eqref{eq:traces}.
\end{itemize}
Permutation tuples are lexicographically ordered. Subset masks use bit $i-1$
for physical site $i$. The checker reconstructs \eqref{eq:Kw} and the
symmetrized auxiliary tensor from \eqref{eq:aweights}, checks
\eqref{eq:cert} coefficient by coefficient, validates every ideal term and
coefficient symmetry, and verifies all polynomial identities and positivity
conditions above. It checks 16,128 polynomial diagram products in the
one-site blocks and all 238 strict-margin conditions \eqref{eq:margin}.

All final positivity tests are on symmetric integer matrices obtained by
clearing positive rational denominators. They pass two exact tests:
fraction-free symmetric elimination and signs of characteristic-polynomial
coefficients. The latter is valid because, for a real symmetric $R$,
nonnegative coefficients of $\det(tI+R)$ exclude a positive root in $t$,
and hence exclude any negative eigenvalue of $R$. Positive Gram limits are
verified separately. No floating-point eigenvalue or optimizer status is
used by the checker.

From the extracted package directory, a complete replay is:
\begin{verbatim}
python -m pip install -r certificate/requirements.txt
python -B certificate/verify.py
\end{verbatim}
The checker resolves its two data files relative to itself. The recorded
replay passes the coefficient and polynomial checks and all 238 blocks by
both exact methods. Its analytic input is precisely
Lemma~\ref{lem:aux}, proved below; the determinant and dimension arguments
have been proved above. These together establish Theorem~\ref{thm:main}.

The SHA-256 digests, fixing the exact electronic data, are
\begin{center}\small\ttfamily
\begin{tabular}{@{}l@{}}
certificate.json.gz:\\
f00440126b5797ca03acc784411873cd15a39652b7dc9674dbb0a1afd759213c\\[3pt]
representations.json.gz:\\
e7652dfef493f38d6ae176e62281b869e310e667480e9c0520cc1ef617f8232f
\end{tabular}
\end{center}
The rational value $\al$ is a certified parameter, not an optimality claim.
The strict margin in \eqref{eq:margin} is a block-coordinate margin, not a
dimension-independent lower bound on $Q_{\al}(X)/\|X\|_F^2$.

\appendix
\section{Proof of the analytic rank-two inequality}\label{app:analytic}
\setcounter{equation}{0}
\renewcommand{\theequation}{A\arabic{equation}}
\begin{lemma}[A rank-two three-factor inequality]\label{lem:aux}
Let $X$ act on a finite-dimensional complex tensor product
$\mathcal H_1\otimes\mathcal H_2\otimes\mathcal H_3$ and have rank at most two.
All partial traces are unnormalized. Put
\[
 N=\|X\|_F^2,\quad T=|\operatorname{Tr}X|^2,\quad
 B_j=\sum_{|E|=j}\|\operatorname{Tr}_E X\|_F^2,\quad
 g=(\sigma_1-\sigma_2)^2,
\]
where a missing singular value is zero. Then
\[
 \frac{17}{8}N-B_1+\frac12B_2-\frac5{16}T\ge\frac g4.
\]
\end{lemma}

\begin{proof}
Write
\[
q(Y)=\sum_{E\subseteq\{1,2,3\}}(-1/2)^{|E|}
                 \|\operatorname{Tr}_E Y\|_F^2.
\]
For every square matrix, $\operatorname{Re}Y=(Y+Y^\dagger)/2$ and
$\operatorname{Im}Y=(Y-Y^\dagger)/(2i)$.
We first prove a positive-matrix bound, then use four positive secants and
a real-Hilbert-space estimate. All identities below remain valid when a
singular value vanishes; the zero matrix is immediate.

\smallskip\noindent\emph{Positive rank-two matrices.}
For $P=W\tau W^\dagger\succeq0$ of trace one, take $W^\dagger W=I_2$ and
$\tau=\operatorname{diag}(\lambda,\mu)$, $\lambda+\mu=1$.
Let $\varsigma_1,\varsigma_2,\varsigma_3$ be the Pauli matrices,
$K_\nu=\sqrt\tau\varsigma_\nu\sqrt\tau$, and
$M_{i\nu}=\operatorname{Tr}_{\widehat i}(WK_\nu W^\dagger)$, where
$\widehat i=\{1,2,3\}\setminus\{i\}$.
In the unit vector $\sum_a\sqrt{\tau_{aa}}|a\rangle W|a\rangle$, the expectation
of $\varsigma_\nu^T\otimes\operatorname{sign}(M_{i\nu})$ is
$\|M_{i\nu}\|_1$. Partition the nine pairs $(i,\nu)$ into the three triples
with $\nu-i$ fixed modulo three. In each triple these Hermitian contractions
anticommute. Thus their squared expectations sum to at most one:
indeed $(\sum_jt_jL_j)^2\preceq\sum_jt_j^2 I$ for anticommuting contractions.
Consequently $\sum_{i,\nu}\|M_{i\nu}\|_1^2\le3$.
For Hermitian $M$, its positive and negative parts give
$2\|M\|_F^2\le\|M\|_1^2+(\operatorname{Tr}M)^2$.
Since $\operatorname{Tr}K_\nu$ is zero except for the value $\lambda-\mu$,
\[
 \sum_{i,\nu}\|M_{i\nu}\|_F^2
       \le\tfrac32[1+(\lambda-\mu)^2]=3\operatorname{Tr}\tau^2.
\]
Pauli completeness
$\sum_{\nu=0}^3(\varsigma_\nu)_{ab}(\varsigma_\nu)_{cd}
=2\delta_{ad}\delta_{bc}$, with $\varsigma_0=I$, gives
\[
 \sum_{\nu=1}^3\|M_{i\nu}\|_F^2
 =2\|\operatorname{Tr}_iP\|_F^2
     -\|\operatorname{Tr}_{\widehat i}P\|_F^2.
\]
Substitution in $q(P)$ and homogeneity therefore prove
\begin{equation}
 q(P)\ge\frac{2\|P\|_F^2-(\operatorname{Tr}P)^2}{8}
       =\frac{\operatorname{gap}(P)^2}{8}.
 \label{eq:psdfloor}
\end{equation}
Here $\operatorname{gap}(P)$ is the difference of its two eigenvalues,
padding by zero.

\smallskip\noindent\emph{Four positive secants.}
Take $X=USV^\dagger$, $U^\dagger U=V^\dagger V=I_2$, with $S$ diagonal
and nonnegative. Set
\[
 s=\operatorname{Tr}S,\quad P=USU^\dagger,\quad Q=VSV^\dagger,
 \quad f=\operatorname{Tr}(PQ),\quad G=2f-T;
 \qquad g=2N-s^2.
\]
Let $F_i$ swap receiver $i$ between two physical replicas,
$\Pi=\prod_i\frac{I-F_i}{2}$, and let $F_\ell$ swap two logical qubits. Define
\[
 \Delta=\sum_{i,\nu=1}^3
 \left\|\operatorname{Tr}_{\widehat i}
\bigl[U\sqrt S\varsigma_\nu\sqrt S U^\dagger
       -V\sqrt S\varsigma_\nu\sqrt S V^\dagger\bigr]\right\|_F^2,
\]
\[
 L=\Pi(U\sqrt S\otimes V\sqrt S),\qquad
 \mathcal S=\operatorname{Tr}\!\left[\frac{I+F_\ell}{2}L^\dagger L\right].
\]
Both quantities are nonnegative. For $z\in\{1,-1,i,-i\}$ put
$P_z=(U+zV)S(U+zV)^\dagger$ and
$e(P_z)=q(P_z)-\operatorname{gap}(P_z)^2/8\ge0$ by \eqref{eq:psdfloor}.
The following exact identity is useful:
\begin{equation}
 24q(X)-4g+G=\sum_{z\in\{1,-1,i,-i\}}e(P_z)+\Delta+64\mathcal S.
 \label{eq:secantidentity}
\end{equation}
For completeness, it can be checked with three short expansions.
Put $E=q(P)+q(Q)$ and define
\[
 h=\sum_J(-1/2)^{|J|}\operatorname{Tr}[(\operatorname{Tr}_JP)
 (\operatorname{Tr}_JQ)],\qquad
 c_j=\sum_{|J|=j}\operatorname{Tr}(P_JQ_J),
\]
where
$P_J=\operatorname{Tr}_{J^c}P$ and similarly for $Q_J$. Pauli completeness
and expansion of the two projectors give, respectively,
\[
 \begin{aligned}
 8h&=8f-4c_2+2c_1-s^2,\\
 4E+\Delta&=8N-s^2-4B_1+2c_1,\\
 64\mathcal S&=4(s^2+N-c_1+c_2-f-B_1+B_2-T).
 \end{aligned}
\]
For the middle equality the cross Pauli sum at receiver $i$ is
$2\|\operatorname{Tr}_iX\|_F^2-\operatorname{Tr}(P_iQ_i)$.
For the last equality the identity and logical-swap contributions are
$(s^2-c_1+c_2-f)/8$ and $(N-B_1+B_2-T)/8$.
Finally, four-phase averaging gives
$\sum_z q(P_z)=4E+8h+8q(X)$ and
$\sum_z\operatorname{gap}(P_z)^2=16g+8G$.
The latter follows by applying
$\operatorname{gap}(M)^2=2\operatorname{Tr}M^2-(\operatorname{Tr}M)^2$
to the two-by-two Gram matrices
$2S+zZ+\bar zZ^\dagger$, $Z=\sqrt S U^\dagger V\sqrt S$.
These expansions prove \eqref{eq:secantidentity}, hence
\begin{equation}
 q(X)\ge(4g-G)/24.
 \label{eq:secantbound}
\end{equation}

\smallskip\noindent\emph{Hermitian reduction and real Bessel inequality.}
For positive $A,B$ with $AB=0$, positivity of
$\prod_i(I-F_i)$ implies $b\le u+a$, where
$u=\operatorname{Tr}A\operatorname{Tr}B$,
$a=\sum_i\operatorname{Tr}[(\operatorname{Tr}_iA)(\operatorname{Tr}_iB)]$,
and $b$ is the analogous sum over two traced factors.
Thus their polarized $q$ value is
$-a/2+b/4-u/8\le u/8$.
Apply this and \eqref{eq:psdfloor} to the positive and negative parts of a Hermitian $Y$
with at most two eigenvalues of either sign, obtaining
\[
 q(Y)\ge[2\|Y\|_F^2-\|Y\|_1^2]/8.
\]
Both $\operatorname{Re}X$ and $\operatorname{Im}X$ have this inertia bound:
their quadratic forms vanish on $\ker X$, a subspace of codimension at most two.
Since $q(X)=q(\operatorname{Re}X)+q(\operatorname{Im}X)$,
\begin{equation}
 q(X)\ge\frac{2N-\|\operatorname{Re}X\|_1^2
                       -\|\operatorname{Im}X\|_1^2}{8}.
 \label{eq:hermitianbound}
\end{equation}
Put $C=U^\dagger V$. Choose a Hermitian unitary
$J=\operatorname{sign}(\operatorname{Re}X)$, assigning signs on its kernel,
and a Hermitian Frobenius-orthonormal basis $E_0,\ldots,E_3$ of $M_2$.
For the real inner product $\operatorname{Re}\operatorname{Tr}(A^\dagger B)$,
the five matrices $JU/\sqrt2,iUE_0,\ldots,iUE_3$ are orthonormal.
Indeed their cross inner products vanish because $U^\dagger JU$ and
the $E_\mu$ are Hermitian. Bessel's inequality on $VS$ gives
\[
 \tfrac12\|\operatorname{Re}X\|_1^2+\|\operatorname{Im}(CS)\|_F^2\le N.
\]
Exchange $U,V$ and average. The elementary entrywise identity
\[
 \|\operatorname{Im}(CS)\|_F^2+\|\operatorname{Im}(SC)\|_F^2
 -2\|\operatorname{Im}(\sqrt S C\sqrt S)\|_F^2
 =\tfrac12\sum_{i,j}(S_{ii}-S_{jj})^2|C_{ij}|^2\ge0
\]
yields $\|\operatorname{Re}X\|_1^2\le2N-2\|\operatorname{Im}Z\|_F^2$.
Apply the same argument to $iX$ and add, using $f=\|Z\|_F^2$:
\[
 \|\operatorname{Re}X\|_1^2+\|\operatorname{Im}X\|_1^2\le4N-2f.
\]
Equation \eqref{eq:hermitianbound} now gives $q(X)\ge(G+T-2N)/8$.
Take one quarter of this bound and three quarters of \eqref{eq:secantbound}:
\[
 q(X)\ge-\frac N{16}+\frac T{32}+\frac g8.
\]
Since $2q(X)+N/8-T/16=17N/8-B_1+B_2/2-5T/16$, the claim follows.
If the entire physical space is one-dimensional, it is immediate directly;
otherwise rank-one SVDs can be padded to the two-column isometries used above.
\end{proof}

\end{document}
